\chapter{Methodology and Implementation}
\label{chap:methodology}

\Cref{chap:architecture} fixed what each component is responsible for and
what data it exchanges with its neighbours. This chapter explains how each
component was built: the concrete rules and libraries used, and the reasons
behind the decisions that shaped them.

The chapter follows the order data actually moves through the system. The
working dataset is prepared first. Then come classification and context
extraction, the explanation component, the PyPI-grounded repair component,
and fix application. The orchestrator that drives these components in a
single unattended pass, and that carries a newly exposed error into the
bounded second round, is described next, because its behaviour is what
the evaluation measures. Result logging and knowledge-graph enrichment
come last, because they join the outputs of every earlier component into
one persistent record.

\section{Implementation Overview}
\label{sec:m-overview}

This section gives one condensed view of how the stages connect, so that
the detail in the rest of the chapter does not obscure how the pieces form
a single system.

\begin{verbatim}
Docker-pipeline dependency-error dataset (214 rows)
  -> dataset preparation      scripts/prepare_dependency_dataset.py
  -> context extraction /
     classification           scripts/extract_error_contexts.py
  -> LLMExplainer             scripts/run_llm_explainer.py
  -> RAGRepairAgent           scripts/rag_repair_agent.py
  -> FixApplicator            scripts/fix_applicator.py
  -> new error: reclassify, explain,
     repair if eligible (round 2)  scripts/run_pipeline.py
  -> ResultLogger (SQLite)    scripts/result_logger.py
  -> CSV export               scripts/export_repair_attempts_csv.py
  -> RML mapping (Morph-KGC)  repair_attempts.rml.ttl
  -> repair RDF linked to FAIR Jupyter KG notebook IRIs
\end{verbatim}

Each arrow is a data hand-off: a plain file, a database, or a call between
the components' own functions. Every stage can be invoked as its own script
against the previous stage's output, which is how each was developed and
tested. The same stages are also driven in sequence by a single
orchestrating process (\texttt{scripts/run\_pipeline.py}) within one pass,
which is how the evaluation in \Cref{chap:validation} was run.

The sections below follow this order, one component at a time.
\Cref{sec:m-worked-example} then traces one concrete record through every
stage.

\section{Dataset Preparation}
\label{sec:m-dataset}

\Cref{chap:problem-dataset} already gave the resulting numbers for this
stage: 214 dependency-error rows, 200 usable and 14 excluded, split 13/187
between development and evaluation. This section explains how
\texttt{scripts/prepare\_dependency\_dataset.py} produces those numbers, and
why the rules behind them were chosen.

The script reads the Docker pipeline's own \texttt{notebook\_executions}
table, and keeps every row where \texttt{error\_category =
'DEPENDENCY\_ERROR'}. For each kept row, it applies two independent
decisions, deterministically and without any LLM involvement. The
first is a \emph{subtype}: is the module missing entirely, present but
incompatible, a missing system library, or an ambiguous local name? The
second is a \emph{scope status}: is this row inside the pip-only v1
repair scope at all? Both decisions rely only on simple, inspectable
rules over the \texttt{error\_type} and \texttt{error\_message} fields
the Docker pipeline already recorded. For example, a message containing
\texttt{".so"} or "shared object file" is classified as
\texttt{system\_library}. A failing module name drawn from a short list
of known-ambiguous names (\texttt{utils}, \texttt{statistics},
\texttt{config}, \texttt{helpers}, \texttt{src}) is classified as
\texttt{mapping\_unknown}.

Doing this filtering with plain rules, before any language model is
involved, was a deliberate choice. A decision that determines whether a
row is even eligible for automated repair should be explainable and
reproducible on its own terms. An LLM's judgement about scope
eligibility would be harder to audit, and could silently change between
runs of the same model. Keeping this gate deterministic also means
every later component can trust \texttt{scope\_status} directly,
instead of re-deriving eligibility itself.

The dataset is also split into three disjoint groups: 13 rows set aside
for manually designing and sanity-checking the explanation and repair
prompts (\Cref{sec:m-explanation,sec:m-repair}), 187 usable rows
reserved for a later, controlled evaluation run, and the 14 excluded
rows. Separating a small development sample from the evaluation set
matters for the same reason it matters in any machine-learning
setting. A prompt tuned by looking at a specific row, and then tested
again on that same row, would give an optimistic and misleading
impression of how well the prompt generalises.

The script writes four output files: a CSV of the classified rows, a schema
document describing every column, a statistics file with the counts used
throughout this thesis, and a split-manifest file recording exactly which
notebook execution IDs fall into which split, so the same split can be
reproduced later rather than re-randomised.

\section{Context Extraction and Error Classification}
\label{sec:m-classification}

\texttt{scripts/extract\_error\_contexts.py} takes each dataset row
prepared above and adds two things: a more detailed root-cause hint
with a confidence level, and, where possible, the actual source code
around the failure. Every field the dataset-preparation stage already
computed (\texttt{scope\_status}, \texttt{exclusion\_reason},
\texttt{split}) is carried through unchanged. This stage enriches
every row regardless of scope, so an excluded row still receives full
context for its explanation.

\subsection{Root-Cause Hints}

The subtype produced by dataset preparation is refined into one of seven
root-cause hints:

\begin{itemize}
  \item \texttt{direct\_missing\_package} -- the module is genuinely absent;
  \item \texttt{import\_distribution\_name\_mismatch} -- the import name
        differs from its PyPI distribution name, for example
        \texttt{sklearn} versus \texttt{scikit-learn};
  \item \texttt{version\_or\_api\_incompatibility} -- a "cannot import name"
        style failure;
  \item \texttt{transitive\_dependency} -- a generic \texttt{ImportError}
        not matching a more specific pattern;
  \item \texttt{system\_level\_dependency} -- a missing shared library;
  \item \texttt{local\_import\_or\_path\_issue} -- an ambiguous local module
        name; and
  \item \texttt{insufficient\_context} -- not enough information to decide.
\end{itemize}

Each hint also carries a confidence level. A direct pattern match gives
high confidence: a missing-module message, a "cannot import name" message,
or a \texttt{.so} failure. Softer phrasing, such as an "install" or
"required" hint, gives medium. An ambiguous local module, or a message
carrying too little information to decide, gives low.

\subsection{Recovering Source Context}

Three sources of context are used, ordered by reliability. The first is a
legacy execution table from an earlier pipeline run. Because this table
was not produced by the same execution as the current dataset and has no
timestamp, it is treated only as supporting evidence.

A legacy row is accepted only when its final reported error matches the
current failure. Matching only the notebook or the failing module name is
not enough. For example, a cell may import both \texttt{pandas} and
\texttt{sklearn}, but fail on \texttt{pandas} in one execution and on
\texttt{sklearn} in another. Treating those as the same failure would
attach the wrong source context.

In practice, 162 of the 214 dataset rows (76\%) have a matching notebook
ID in the legacy table, but only about 40\% of those records contain the
correct failing module in the stored traceback. For this reason, legacy
context is used only as a low-confidence hint and is never trusted on its
own.

The second tier, used only when explicitly requested, fetches the
notebook and its dependency files directly from GitHub. It uses the
repository's recorded commit where available, and falls back to the
\texttt{main} or \texttt{master} branch otherwise. The third tier
falls back to the metadata already present from dataset preparation
(error type, message, cell index, failing module), when no source can
be recovered at all. This is unavoidable here: the repositories cloned
during the original Docker pipeline run were never kept on this
machine, so no local copy of any notebook's source exists unless it is
fetched again.

\subsection{Limitations}

A few limitations follow directly from the above, and are worth stating
plainly, because they affect how confidently later components can rely
on this context. Remote fetching is not pinned to the exact commit that
originally failed, unless that commit happens to still be resolvable.
So a freshly fetched notebook can, in rare cases, differ slightly from
the one that actually produced the recorded error.
\texttt{error\_cell\_index} indexes into the notebook's full cell list,
so a notebook edited since the original run can shift cell positions
and misalign the index. Any unexpected per-row failure (a malformed
row, a database error) falls back to a bare metadata-only record,
rather than aborting the whole batch -- this keeps one bad row from
blocking the other 213.

\section{LLM Explanation Component}
\label{sec:m-explanation}

LLMExplainer (\Cref{subsec:llmexplainer}) turns a classified failure into a
plain-language account of what went wrong. It is the component most
directly tied to the thesis title, and its role is deliberately narrow: it
explains a failure, it never proposes how to fix one.

\subsection{Role and Separation from Repair}
\label{sec:m-explanation-role}

The prompt sent to the model instructs it to describe the failure in plain,
non-technical language, using only the information supplied, and explicitly
forbids it from suggesting a fix -- deriving a repair is RAGRepairAgent's
job, kept structurally separate (\Cref{sec:m-repair}). This separation
serves two purposes. First, it keeps each component's failure mode
contained: an explanation that misjudges a detail is a readability problem,
while a repair suggestion that misjudges a package name or version is
executed inside a container, so the two carry very different risk and
deserve very different validation. Second, it lets LLMExplainer's scope be
wider than RAGRepairAgent's: the pip-only repair scope defined in
\Cref{sec:pd-scope} restricts which rows may be \emph{repaired}, not which
rows may be \emph{explained} (\Cref{sec:pd-explanation-scope}). A researcher
whose notebook fails because of a missing system library still receives a
plain-language explanation of the failure, even though the current repair
layer cannot act on that particular case.

The same two principles carry over to the second repair round. When a
repaired notebook fails again on a new error, that error is reclassified
and explained as a record of its own before its repair eligibility is
evaluated, so the reader receives an account of every error the pipeline
encountered, not only the first (\Cref{sec:m-orchestrator-round2}). The
explainer is called with the same prompt, schema, and retry rule in both
rounds; the round only changes which record it is given.

\subsection{Model and Runtime Selection}

The explanation and repair components share one open, locally-run
model, rather than each using a different one. \textbf{Gemma-2 9B} was
selected as the primary model because PLLM~\cite{bartlett2025pllm}, the
closest prior system, reports Gemma-2 9B with RAG as its own
best-performing configuration for exactly this kind of dependency task.
Using the same model for both explanation and repair keeps a later
comparison between configurations clean: any difference in outcome can
be attributed to the prompt or the retrieval design, not to a hidden
difference in which model produced which part of the result.
CodeLlama-13B and DeepSeek-Coder (7B and 13B) were kept as fallback
candidates, motivated by Szalontai et al.~\cite{szalontai2024codellama}
and a broader empirical comparison of code-specialised
models~\cite{empiricalapr2025}. Both report that smaller,
code-specialised models can match or beat larger general-purpose ones on
repair tasks. Neither fallback was needed: Gemma-2 9B performed adequately
in the sanity checks run on the development split. A separate experiment
later replaced Gemma with a much larger open model to test how much the
choice mattered (\Cref{sec:v-model-sensitivity}).

Ollama was chosen as the serving runtime because it gives simple local
access to a quantised model without setting up a GPU-serving stack, which
is appropriate for a component running one notebook at a time rather than
serving high-throughput traffic. A hardware check confirmed \texttt{gemma2:9b}
runs on the available machine using 6.3~GB of memory, entirely on CPU, with
a 4096-token context window -- workable, though slow, which later became
relevant to a timeout decision described in \Cref{sec:m-repair}.

\subsection{Structured Explanation Output}
\label{sec:m-explanation-deviation}

An explanation is read by a person, but it also has to be logged,
validated, and processed by later stages. Free prose cannot be checked
automatically for the fields those stages need. The explainer therefore
returns a single JSON object, validated against a fixed schema
(\texttt{schemas/}\allowbreak\texttt{explanation.schema.json}) with exactly
six fields:

\begin{itemize}
  \item \texttt{summary} -- a short, plain-language statement of what failed;
  \item \texttt{root\_cause} -- why it failed, in the same plain language;
  \item \texttt{evidence} -- a list of the concrete signals the explanation
        is grounded in (the error type, the error message, the root-cause
        hint);
  \item \texttt{failing\_module} -- the module named in the explanation;
  \item \texttt{explanation\_confidence} -- \texttt{high}, \texttt{medium},
        or \texttt{low}; and
  \item \texttt{limitations} -- what the explanation could not determine,
        for example because only metadata, not source code, was available.
\end{itemize}

The plain-language requirement is enforced inside the fields rather than
abandoned: \texttt{summary} and \texttt{root\_cause} are what a reader
sees, and the prompt still constrains them to short, non-technical
sentences. The surrounding structure exists for the pipeline, not for the
reader. This follows the same pattern used for the repair prompt
(\Cref{sec:m-repair}).

The schema's \texttt{explanation\_confidence} field is named differently
from the classifier's own \texttt{confidence} field
(\Cref{sec:m-classification}) because the two measure different things and
can legitimately disagree. A classifier can be highly confident it matched
a known pattern while the model is only moderately confident it has
explained the underlying cause correctly.

\subsection{Prompt, Validation, and Retry}

The prompt instructs the model to write in plain, non-technical language,
to use only the information supplied, and never to suggest a fix. Using
only supplied information means never inventing a version number, a file
name, or a cause not visible in the error. The runner sends
this prompt to \texttt{gemma2:9b} through Ollama (or, in the
model-sensitivity run of \Cref{sec:v-model-sensitivity}, to Qwen through
the same provider interface with the same prompt), parses the response as
JSON, and validates it against the schema. Each request has a 120-second
timeout. On invalid JSON, a schema violation, a timeout, or an unavailable
model, the runner retries once with a corrective prompt. If the second
attempt also fails, the failure itself is logged, with its category
(\texttt{timeout}, \texttt{model\_unavailable}, or a validation failure),
rather than discarded, so no row silently disappears from the results
file. The evaluation in \Cref{chap:validation} keeps these categories
apart: a request that timed out is a runtime failure of the model service,
not a malformed model output, and the two are reported separately.

Each logged result keeps three things apart: the input metadata from the
classification stage, the model and prompt configuration used, and the
generated explanation with its validation status, latency, and retry count.
Separating them means a later reader cannot confuse the original error text
with text the model wrote about it. It also lets the result-logging stage
(\Cref{sec:m-resultlogger}) join an explanation back to its originating
notebook by \texttt{notebook\_execution\_id} alone.

\subsection{A Real Explanation Record}
\label{sec:m-explanation-example}

The following committed record illustrates the schema above on a real,
already-classified failure: \texttt{notebook\_execution\_id}~8, a
\texttt{ModuleNotFoundError} for \texttt{sklearn}.

\begin{quote}
\small
\begin{sloppypar}
\textbf{summary:} "The notebook execution failed because Python could not
find the module 'sklearn'."\\[2pt]
\textbf{root\_cause:} "The notebook imports 'sklearn', which is not
available in the execution environment. The classifier also marks this as
an import/distribution-name mismatch, meaning the import name may not
directly match the package distribution name."\\[2pt]
\textbf{evidence:} ["The error type is ModuleNotFoundError.", "The error
message says: No module named 'sklearn'.", "The root\_cause\_hint is
import\_\allowbreak distribution\_\allowbreak name\_\allowbreak mismatch."]\\[2pt]
\textbf{failing\_module:} "sklearn"\\[2pt]
\textbf{explanation\_confidence:} "high"\\[2pt]
\textbf{limitations:} "Only metadata is available, so the exact failing
cell and dependency files cannot be inspected."
\end{sloppypar}
\end{quote}

This record was generated with the \texttt{few\_shot} prompt strategy. It
succeeded on the first attempt and required 98{,}029~ms on the CPU-only
hardware described above. The example is included only to illustrate the
structure of a real explanation record; it should not be interpreted as
evidence of explanation quality across the full dataset.
The few-shot prompt contains two fixed examples in\\
\texttt{prompts/dependency\_explanation\_v1.txt}:
this \texttt{sklearn}/\texttt{ModuleNotFoundError} case and the
\texttt{scipy}/\texttt{cumtrapz} case discussed in
\Cref{sec:m-worked-example}. Both examples come from the 13-row development
split and remain unchanged for every LLM call. No example from the reserved
evaluation split is included in the prompt.


\section{PyPI-Grounded Repair Generation}
\label{sec:m-repair}

RAGRepairAgent (\Cref{subsec:repairagent}) does not install anything. Its
only job is to answer one question -- \emph{what repair should be
attempted?} -- and to answer it safely: never inventing a package name,
never inventing a version number, and never handing a later component a raw
string to execute in a shell. This section explains how that is achieved.

\subsection{Overall Flow}

\begin{verbatim}
classified dependency failure (usable)
    -> identify the failing import
    -> resolve the import name to a PyPI distribution name
    -> query the PyPI registry for that distribution
    -> filter the returned versions to a small, safe candidate set
    -> for a version-incompatibility failure, intersect with
       curated compatibility evidence
    -> give the LLM only this grounded candidate set
    -> validate the LLM's proposal against the same evidence
    -> if valid, deterministically build the exact install command
\end{verbatim}

The design principle behind this flow is that the language model is never
the only thing standing between an error message and a command that will
later run inside a container. Four separate layers each have exactly one
job, and each layer only trusts what the previous layer has already
verified.

\subsection{Layer 1: Deterministic Retrieval}
\label{sec:m-repair-retrieval}

Before any network request, the failing Python import name must be
resolved to an installable PyPI distribution name. The two are
frequently different: \texttt{sklearn} installs as
\texttt{scikit-learn}, \texttt{cv2} as \texttt{opencv-python},
\texttt{skimage} as \texttt{scikit-image}, \texttt{umap} as
\texttt{umap-learn}, \texttt{pkg\_resources} as \texttt{setuptools}, and
\texttt{Bio} as \texttt{biopython}. A small, explicit, hand-curated
mapping table (\texttt{config/package\_mapping.yaml}) holds every case
the retriever is allowed to resolve. An import name absent from this
table is never assumed to also be its own package name. Instead, the
retriever reports a \texttt{mapping\_unknown} status, without making
any PyPI request at all. This rule matters concretely for this
thesis's own dataset: \texttt{sklearn} and \texttt{pkg\_resources} are,
respectively, the second and first most common failing modules
(\Cref{tab:pd-top-modules}). Getting this mapping right or wrong
affects a large share of the dataset, not an edge case.

Once resolved, the distribution's release metadata is retrieved from the
official PyPI JSON Simple API
(\texttt{GET /simple/\{distribution\}/}), normalising the name first
according to PEP 503. This endpoint was chosen deliberately over the older
project JSON endpoint's \texttt{releases} field, which PyPI's own
documentation marks as deprecated for new integrations; the newer endpoint
is the one PyPI recommends going forward. The returned files are grouped by
release version, and a safe candidate set is computed by dropping
unparseable files, pre-releases, and releases where every file is yanked,
and by checking each release's declared \texttt{requires-python} against
the runtime the notebook actually executed under.

That runtime version is a fixed constant, \texttt{"3.10"}
(\texttt{config/rag\_repair.yaml}), not looked up per notebook. This is not
an approximation of convenience: the Docker pipeline builds every
repository's container from the same template, with the base image
hardcoded to \texttt{python:3.10-slim}, so every notebook this component
processes genuinely ran under the same interpreter version. The
notebook's own recorded kernel metadata is deliberately not used as a
substitute, because it reflects what the notebook's original author had
installed, which is frequently stale or missing, not what the Docker
container the notebook actually failed in provides.

The final candidate list is capped to five versions, newest first, to keep
the amount of version metadata placed in the prompt bounded and readable.

\subsection{Why PyPI Metadata Is Not Enough}
\label{sec:m-repair-compat}

For a \texttt{wrong\_version} failure -- a "cannot import name" error --
knowing that a release exists and supports the right Python version is
not enough. PyPI's own Simple API specification (PEP 691) only
exposes, per file, its filename, URL, hashes, required Python version,
size, upload time, and yanked status. It says nothing about which
functions or classes a given release actually provides. Proving that,
for example, \texttt{scipy.integrate.cumtrapz} still exists in a
specific SciPy release requires consulting that project's own
documentation, not PyPI.

This gap is closed with a small, hand-curated compatibility registry. Each
entry records one import symbol, the version range in which an official
source confirms that symbol still exists, and a link to that source. \Cref{tab:m-compat}
lists the three patterns curated for this dataset, all covering symbols
that were removed in a major SciPy or NumPy release.

\begin{table}[htbp]
  \centering
  \begin{tabular}{|p{3.8cm}|p{2.5cm}|p{5.7cm}|}
    \hline
    \textbf{Symbol} & \textbf{Compatible with} & \textbf{Official source} \\
    \hline
    \texttt{scipy.\allowbreak integrate.\allowbreak cumtrapz} & \texttt{<1.14.0} & SciPy 1.14.0 release notes, "Expired deprecations" \\
    \hline
    \texttt{scipy.sparse.\allowbreak sputils.\allowbreak isshape} & \texttt{<1.14.0} & SciPy 1.14.0 release notes, cross-checked against SciPy's source repository \\
    \hline
    \texttt{numpy.\allowbreak Visible\allowbreak Deprecation\allowbreak Warning} & \texttt{<2.0.0} & NumPy 2.0.0 release notes, "NumPy 2.0 Python API removals" \\
    \hline
  \end{tabular}
  \caption{Curated API-compatibility evidence used for \texttt{wrong\_version} repairs.}
  \label{tab:m-compat}
\end{table}

The registry is intentionally small and manually curated. It contains verified
compatibility rules only for the specific wrong-version API failures observed
in this thesis dataset. Each rule was checked against an official source. If a
wrong-version failure is not covered by the registry, the system does not guess
a compatibility rule and therefore abstains from proposing a version-based
repair.
The registry is therefore dataset-specific rather than a general-purpose
compatibility database. Its purpose is to provide reliable evidence for the
wrong-version patterns present in the evaluated corpus.

\subsubsection{Layer 2: What the Model Decides}

After Layer~1 has resolved the package and retrieved the available PyPI
evidence, the repair model chooses a repair from that grounded information.
It receives the classified failure, including the error type and message,
failing module, subtype, and root-cause hint. Where available, the failing
cell source is also included. Layer~1 additionally provides the resolved
distribution name, the Python target, the compatible candidate releases,
and, for wrong-version failures, the applicable compatibility constraint.

The model does not receive the raw PyPI response. Instead, it sees only the
candidate set already filtered by the deterministic retrieval layer. It also
does not receive the LLMExplainer output or repository dependency files such
as \texttt{requirements.txt}; the latter are used later when FixApplicator
reconstructs the execution environment.

The model returns a structured repair proposal with three possible actions:
\texttt{install}, \texttt{pin\_version}, or \texttt{none}. For
\texttt{install}, it identifies the distribution to install. For
\texttt{pin\_version}, it additionally selects a version from the supplied
candidate set. The \texttt{none} action represents abstention when the
available evidence does not support a repair.


\subsection{Layers 3 and 4: Validation and Command Construction}

Before a repair proposal is used, it first passes schema validation and is then 
checked against the evidence retrieved in Layer 1. 
Five deterministic grounding checks are applied:

\begin{sloppypar}
\begin{itemize}

  \item \texttt{action} must be one of the three allowed values:
        \texttt{install}, \texttt{pin\_version}, or \texttt{none}.

  \item A proposed \texttt{install\_name} must exactly match the
        distribution resolved by the retriever.

  \item A proposed \texttt{version} must appear in the supplied
        candidate list.

  \item An \texttt{install} action for a \texttt{missing\_package}
        additionally requires at least one Python-compatible candidate.

  \item A \texttt{pin\_version} action for a \texttt{wrong\_version}
        additionally requires verified compatibility evidence.

\end{itemize}
\end{sloppypar}

The proposal is also checked for unsafe string values, such as whitespace,
leading dashes, or shell metacharacters. If any validation check fails, the
proposal is rejected and the action is set to \texttt{none}. The system does
not replace an invalid package or version with a guessed alternative.

Only validated proposals are converted into executable commands. The command
is constructed deterministically by the pipeline rather than generated by the
LLM. For example, an installation produces
\texttt{python -m pip install <package>}, while a version pin uses
\texttt{<package>==<version>}. This ensures that unrestricted model-generated
text is never executed directly.

If the model response cannot be parsed, fails schema or grounding validation,
times out, or the model is unavailable, the component retries once. The reason
for the failure is included in the retry prompt. If the second attempt also
fails, the component returns \texttt{action: none}, and no repair is applied.

\subsection{Why Repository Dependency Files Are Not Used as Repair Evidence}
\label{sec:m-repair-nodeps}

Repository dependency files such as \texttt{requirements.txt} are useful for
reconstructing an environment, but they are not reliable enough to serve as
direct repair evidence. A declared dependency file does not necessarily
describe the exact package state in which the recorded failure occurred. It
may be incomplete, outdated, unpinned, or may itself contain the version that
caused the incompatibility.

The files are also inconsistent across repositories. Some projects use\\
\texttt{requirements.txt}, while others use \texttt{setup.py},
\texttt{pyproject.toml}, multiple requirement files, or no explicit dependency
file at all. Making repair generation depend on one of these files would
therefore provide different amounts and qualities of evidence for different
notebooks.

For these reasons, the repair decision is based on evidence that can be checked
independently of the repository's declared environment: the observed failure,
the resolved PyPI distribution, available releases, and compatibility evidence
where applicable. Repository dependency files are used later by FixApplicator
when rebuilding the baseline environment in which the proposed repair is
tested.

This separation gives the two kinds of information different roles. Dependency
files describe how the repository expects its environment to be constructed,
whereas the repair evidence is used to determine how the observed dependency
failure should be addressed.

\subsection{Runtime Configuration}

RAGRepairAgent shares its model with the explanation component:
\texttt{gemma2:9b} through Ollama, temperature 0.1, top-p 0.9, and at most
700 generated tokens. The request timeout is 300 seconds. This value was
set from measurement. In a pilot run on the CPU-only hardware described
above, the same repair request timed out twice under an initial 120-second
limit and then succeeded at the first attempt under a 300-second limit, on
the same machine with the same model (\Cref{sec:v-i4}). The model was
healthy throughout; 120 seconds was simply too short for CPU-only
generation.

The retrieval client enforces a 0.5-second wait between PyPI requests and
at most one retry on an HTTP 429 response, so a batch run over many rows
cannot overload the public registry.

\subsection{The Implemented Result Shape}

\begin{sloppypar}
Each persisted result record carries three groups of fields.
\texttt{notebook\_execution\_id} is its only join key back to the rest of
the dataset. An \texttt{eligibility} block records whether the row was
usable, excluded, or invalid. A nested \texttt{retrieval\_result} block
holds everything Layer~1 found. The validated outcome is stored as
\texttt{final\_action}, \texttt{final\_install\_name}, and
\texttt{final\_version}, together with the constructed argument list and
its human-readable \texttt{command} string.
\end{sloppypar}

\begin{sloppypar}
An overall \texttt{status} field summarises the attempt with one of three
values. \texttt{"success"} marks a grounded proposal. \texttt{"abstained"}
marks a legitimate \texttt{none}, whether the model chose it or a failed
check forced it. \texttt{"failed"} means every attempt was exhausted
without a usable response. The evaluation in \Cref{chap:validation} counts
abstentions from this three-way distinction.
\end{sloppypar}

The fix object stores only the repair information and the
\texttt{notebook\_execution\_id} used as a join key. Repository and notebook
metadata are not duplicated in this object. 
Before applying a repair, FixApplicator uses the join key to recover the
corresponding repository and notebook information from the original dataset
(\Cref{sec:m-fixapply}).

\section{Fix Application and Notebook Re-execution}
\label{sec:m-fixapply}

FixApplicator answers the question RAGRepairAgent leaves open: \emph{does
the proposed repair actually work?} It contains no language-model call. It
is a small, deterministic pipeline that applies an already-validated
proposal, re-runs the notebook, and classifies what happened.

\begin{verbatim}
validated RAGRepairAgent proposal
    -> re-join with the classification dataset
       by notebook_execution_id
    -> re-validate install_name/version and rebuild the command
    -> rebuild a Docker environment from the original recipe
    -> apply the fix inside that environment
    -> re-run the notebook top to bottom
    -> classify the outcome
\end{verbatim}

\subsection{Recovering Notebook Identity}

\begin{sloppypar}
A RAGRepairAgent result is joined to the rest of the dataset by
\texttt{notebook\_execution\_id} alone; it carries no repository ID,
notebook name, or repository URL. FixApplicator therefore first looks up
each record's \texttt{notebook\_execution\_id} in the classification
dataset
(\texttt{data/context-classification/\allowbreak dependency\_error\_contexts.jsonl})
to recover this identity. If the ID is not found, or the matching record
lacks one of these fields, FixApplicator records the attempt as an
infrastructure failure (\texttt{apply\_error}, failure stage
\texttt{"join"}). Filling in a missing field with a default value could
silently clone or run the wrong repository, so no guess is made.
\end{sloppypar}

\subsection{Re-validating the Persisted Command}

FixApplicator does not execute the persisted \texttt{argv} as it is, even
though RAGRepairAgent already validated the proposal. It re-checks
\texttt{install\_name} and \texttt{version} with the same safety check
RAGRepairAgent uses, and rebuilds the argument list from scratch with the
same function. If the rebuilt list differs from the persisted one, the
record is rejected before anything runs. This second check enforces the
same safety rule independently. When the two components agree, the
proposal has not been altered between the time it was written and the time
FixApplicator reads it.

\subsection{Rebuilding the Docker Environment}
\label{sec:m-fixapply-docker}

\begin{sloppypar}
The upstream Docker pipeline does not preserve its per-repository
containers, so there is no original container to reuse. FixApplicator
therefore rebuilds the environment for every repair attempt from the same
recipe, as \Cref{chap:architecture} specifies: the same base image
(\texttt{python:3.10-slim}), the same dependency-installation procedure,
and the same notebook execution command
(\texttt{jupyter nbconvert --to notebook --execute --allow-errors}) that
the reference pipeline uses.
\end{sloppypar}

The installation procedure is where the repository's own dependency files
enter the pipeline. A \texttt{requirements.txt} at the repository root,
when present, is installed entry by entry; an entry that fails to install
is skipped rather than aborting the build. Afterwards, each
\texttt{setup.py} whose location the upstream repository metadata records
is installed from its directory. These files define the baseline
environment in which the fix is tested; they are not evidence for the fix
itself (\Cref{sec:m-repair-nodeps}). The \texttt{setup.py} locations come
from the upstream metadata rather than from scanning the clone. In the
final evaluation runs, that metadata lookup returned no data
(\Cref{sec:v-final-manifest}), so the baseline there consisted of the
root-level \texttt{requirements.txt} where the repository had one.

The execution command is reused intentionally instead of writing a new
execution loop: it is the exact mechanism that produced the original
dataset, so reusing it keeps the comparison fair. The rebuilt environment
is meant to keep environmental differences small, not to be an exact copy.
The repository is cloned again, and packages the repository does not pin
resolve to whatever PyPI serves at run time. Only the one notebook that
failed is executed, not every notebook in the repository, and it is always
executed from the beginning. A notebook keeps state across cells, so a
partial re-run from the failing cell cannot show that a fix works.

Two implementation details differ from the reference Docker pipeline. First, 
FixApplicator keeps the Docker build context separate from the cloned
repository. This prevents the repository's full Git history from being sent to
the Docker daemon during the build. The repository is still mounted into the
container when the notebook is executed, so this does not change the notebook
source available at run time. Second, the temporary repair container runs as the 
default root user instead of being mapped to the host user's ID. This keeps the 
execution setup simple and allows the already validated installation command to 
be executed without modification. The container is disposable and is used only 
for validating the repair.

\subsection{Commit Pinning}

Where the upstream pipeline's own database records a commit for a
repository, FixApplicator checks out that exact commit before building
anything, rather than staying on the repository's current default branch.
A public GitHub repository can change after the original dataset was
collected; without pinning, a later fix attempt could be tested against a
different version of the notebook than the one that produced the original
failure. The reference pipeline does not pin the commit itself: it clones
only the latest state of the default branch, even though it records a
commit. Here the implementation goes further than the system it builds
on. If the recorded commit cannot be checked out, FixApplicator reports a
specific infrastructure failure rather than falling back silently to the
current branch.

The commit lookup is optional. If the metadata database is unavailable,
FixApplicator continues without commit pinning and records that the checkout
was skipped. In that case, the repository is cloned from its current default
branch. This fallback was used in the final evaluation runs because the metadata
database configured for FixApplicator was not accessible from the WSL
environment. As a result, the evaluated repositories were not pinned to their
original recorded commits. \Cref{sec:v-final-manifest}  discusses this
limitation and its effect on reproducibility.

\subsection{Fix Application Safety}

The install command is executed only as an explicit list of arguments
passed directly to a subprocess, never as a shell string. Every external
process this component invokes, git and Docker included, goes through an
injectable runner as an argument list. A static check of the source code
confirms that \texttt{shell=True}, \texttt{os.system}, \texttt{eval}, and
\texttt{exec} appear nowhere in it. If the fix installation fails inside
the container, the container script stops immediately and the notebook is
not executed, because running an unfixed notebook would validate nothing.

\subsection{Outcome Classification}

Exactly three outcomes are possible, listed in \Cref{tab:m-outcomes}. They
match the contract fixed in \Cref{chap:architecture}; no fourth value is
introduced.

\begin{table}[htbp]
  \centering
  \begin{tabular}{|l|p{9.3cm}|}
    \hline
    \textbf{Outcome} & \textbf{Meaning} \\
    \hline
    \texttt{fixed} & The re-executed notebook produces no error output anywhere. \\
    \hline
    \texttt{still\_failing} & The fix installed and the notebook re-ran, but an error remains -- possibly the same error as before, possibly a different one further into the notebook. \\
    \hline
    \texttt{apply\_error} & The repair could not be meaningfully tested: a dataset join failure, an invalid repair-proposal record, a clone or checkout failure, a Docker build failure, a failed fix installation, a timeout, or a missing output notebook. \\
    \hline
  \end{tabular}
  \caption{FixApplicator outcome values.}
  \label{tab:m-outcomes}
\end{table}

A \texttt{still\_failing} result additionally records whether the new
error is the same as the original one. The check is conservative: the
error is marked as the same only when the exception type matches
\emph{and} the original failing module name still appears in the new
message. The check is designed to under-claim sameness rather than
over-claim it, so a fix that resolved the original problem and uncovered
a different one is not reported as having done nothing. This distinction
matters a great deal in practice, as \Cref{chap:validation} shows.

\begin{sloppypar}
A row whose proposal was \texttt{"none"}, or whose upstream status was not
\texttt{"success"}, is not treated as an attempt. It is recorded with a
separate \texttt{status: "skipped"} field rather than a fourth outcome
value. This keeps the three-value contract intact while still separating
``nothing was attempted'' from ``an attempt was made and failed''.
\end{sloppypar}

\subsection{Why a Real Pilot Was Necessary}
\label{sec:m-fixapply-bugs}

The FixApplicator test suite passed before the component was tested with real
Docker executions. However, the first real pilot revealed three
environment-specific problems. Docker output was decoded incorrectly, the
generated entry-point script had incompatible line endings, and cleanup failed
for some read-only files created by \texttt{git clone} on Windows. These
problems were fixed, and regression tests were added for each of them.

This showed why real execution tests were necessary. Mocked tests can verify
the internal logic of FixApplicator, but they do not reproduce all behaviour
of Docker, Git, the operating system, and subprocess execution. Therefore, the
pilots in \Cref{chap:validation} use real executions in addition to the automated test
suite.


\subsection{Environment Isolation and What Is Out of Scope}

Every attempt uses a freshly named work directory and container. Both are
removed afterwards, whether the attempt succeeded, failed, or timed out.
Docker may reuse a built image across attempts through its own build
cache; this is expected and does not break isolation. The guarantee that
matters is a fresh clone and a fresh container per attempt, not a fresh
image layer every time.

FixApplicator returns one self-contained result per attempt. It has no
orchestrator and no two-round loop, and it does not persist its results to
a database or knowledge graph. ResultLogger persists the results
(\Cref{sec:m-resultlogger}), and the orchestrator, described next, carries
a newly exposed error into the bounded second round.

\section{Orchestration and the Bounded Second Round}
\label{sec:m-orchestrator-round2}

The orchestrator, \texttt{scripts/run\_pipeline.py}, drives the components
above over a dataset split in one unattended pass, calling each through
its own interface. Its control flow was fixed in \Cref{sec:orchestrator}.
This section describes how that flow is realised, in particular how a
newly exposed error is handled, because the evaluation in
\Cref{chap:validation} measures exactly this behaviour.

\subsection{One Pass per Notebook}

For each record of the split, the orchestrator classifies the record and
calls LLMExplainer on it. If the record is repair-eligible, it calls
RAGRepairAgent, and if the agent produces a grounded proposal, it calls
FixApplicator. Every call for one record shares one run identifier. Round
1 is the same for every record. A second round is considered only when
three conditions hold: Round 1 reached FixApplicator, the notebook came
back \texttt{still\_failing}, and the conservative same-error check
(\Cref{sec:m-fixapply}) reports a genuinely different error. A notebook
that still fails on its original error, or whose repair could not be
tested, receives no second round, because repeating the same repair would
test nothing new.

\subsection{Reclassifying and Explaining a Newly Exposed Error}
\label{sec:m-orchestrator-round2-explain}

A genuinely new error is turned into a record of its own. The orchestrator
runs the same two classification stages that every original row went
through (\Cref{sec:m-dataset,sec:m-classification}) on the new error type
and message. The result is a \emph{round-2 record} with its own failing
module, subtype, scope status, and root-cause hint. The record is
explained first. A helper, \texttt{explain\_round\_record()}, calls the
same explainer on the round-2 record with the same prompt template,
schema, and one-retry rule as Round 1, and tags the resulting explanation
with its round. Two properties of this step matter for the evaluation.

First, the explanation is produced for every newly exposed error,
whatever its scope. A new \texttt{missing\_package} error is explained and
then repaired. A new \texttt{system\_library} error, a new
\texttt{mapping\_unknown} error, or a non-dependency error is explained,
and the record then stops there because it is not repair-eligible.
Explanation scope therefore stays wider than repair scope in Round 2,
exactly as it does in Round 1.

Second, the explanation is independent of the repair. RAGRepairAgent
receives the classified round-2 record and the PyPI evidence retrieved for
it, never the explanation text. A failed explanation call does not block
repair either. If the explainer times out, the model is unavailable, or
the call raises unexpectedly, the helper records a failed explanation with
its failure category and returns normally. An otherwise eligible record
then proceeds to RAGRepairAgent and FixApplicator as usual, and the
evaluation counts it as a repair round with a failed explanation.

Repair eligibility is evaluated only after the explanation is recorded,
with the same rules as in Round 1. An eligible round-2 record is passed to
RAGRepairAgent as a target in its own right, with its own retrieval,
proposal, and validation. Two rounds is a hard cap
(\texttt{MAX\_ROUNDS\_HARD\_CAP = 2}). An error exposed by the Round-2
re-execution is recorded in the Round-2 outcome but is not reclassified,
explained, or repaired. No record therefore carries more than two
explanations or two repair rounds.

\subsection{Why Round 2 Reapplies Round 1's Fix}
\label{sec:m-fixapply-round2}

The isolation guarantee of \Cref{sec:m-fixapply} has a direct consequence
for the second round: no container survives between Round 1 and Round 2 of
the same notebook. If Round 2 rebuilt a fresh container and applied only
its own fix, the notebook would be tested against \emph{original
environment + Round 2}, without Round 1's already-validated fix. A
\texttt{fixed} or \texttt{still\_failing} verdict from that container
would show only whether the second repair alone was enough. The question
Round 2 is meant to answer is whether both repairs together make the
notebook run.

The orchestrator therefore re-runs Round 1's already-validated install
command inside Round 2's freshly rebuilt container, immediately before it
applies Round 2's own fix. The notebook is re-executed against
\emph{original environment + Round 1 + Round 2}. Every Round-2 outcome
then reflects what it claims to validate, and no container has to persist
between rounds: the fresh-container-per-attempt guarantee still holds for
every Round-2 attempt, and only the fixes applied inside that fresh
container change.

\subsection{What Is Recorded}
\label{sec:m-orchestrator-record}

The orchestrator writes two things. The first is a per-notebook JSONL
trace: one line per notebook (187 in the evaluation split), holding the
classification, the Round-1 explanation, the second-round decision, and
one entry per repair round in which RAGRepairAgent was invoked. A
reclassified new error's explanation is stored with the decision that
produced it, so it is present whether or not a second round followed. The
top-level explanation of a trace line is always the original failure's,
so a trace from a single-round run has the same fields in the same place.

The second is the \texttt{repair\_attempts} table, written through
ResultLogger (\Cref{sec:m-resultlogger}) with one row per repair-decision
entry: normally a RAGRepairAgent invocation, and for an original
out-of-scope record a synthesised abstention entry written without an
agent call (\Cref{sec:m-resultlogger-join}). An abstention is a row too,
with \texttt{action = none} and no
outcome. A Round-1 row carries the explanation, model, and prompt strategy
of the original failure; a Round-2 row carries those of the newly exposed
error it targeted. A newly exposed error that was explained but was not
repair-eligible never reaches RAGRepairAgent and therefore has no Round-2
row: its explanation exists in the trace only. In the final evaluation the
table holds 235 rows per run, 187 for Round 1 and 48 for Round 2, of which
73 carry an executed fix application. The table's schema is unchanged by
any of this: the same eighteen columns hold either round's values.

\section{Result Logging and Knowledge-Graph Enrichment}
\label{sec:m-resultlogger}

ResultLogger (\Cref{subsec:resultlogger}) classifies, explains, proposes,
and applies nothing. It joins the JSONL output files of the preceding
components, one file per component, into one persistent
\texttt{repair\_attempts} row per repair-decision entry, and makes that
row exportable for the knowledge-graph mapping. Three files implement it.
\path{scripts/result_logger.py} covers the SQLite side.
\path{scripts/export_repair_attempts_csv.py} and
\path{mapping/rml_mapping/repair_attempts.rml.ttl} cover the
knowledge-graph side.

\subsection{Joining Four Outputs}
\label{sec:m-resultlogger-join}

RAGRepairAgent's repair-proposal output is the driving input: ResultLogger
writes exactly one row per entry in it, in order. An invocation that
abstains is an entry too, so the number of entries is the number of repair
decisions taken, not the number of classified or explained errors. Repair
eligibility is checked before RAGRepairAgent is called
(\Cref{sec:m-orchestrator-round2}), and the two rounds then differ in what
an ineligible record produces. For an original record outside repair
scope, the orchestrator synthesises an entry in the same abstained result
shape without any model call, reusing the agent's own eligibility check
and result skeleton (\texttt{build\_excluded\_repair\_stub} in
\path{scripts/run_pipeline.py}), so that every original record still
contributes one row. A newly exposed error is instead reclassified and
explained, and reaches RAGRepairAgent only if it is itself repair-eligible;
a non-repairable new error stops there, its explanation existing in the
trace only and producing no row. In the final evaluation every one of the
187 original records was repair-eligible, so no synthesised entry occurs
in it: 50 newly exposed errors were explained, 48 became Round-2
invocations, and all 187 + 48 = 235 entries are genuine invocations,
giving 235 rows.

The other three inputs are joined to that driving record in three
different ways, because the data does not offer one uniform key. The
classification dataset is joined by \texttt{notebook\_execution\_id}; a
miss is a hard error, as it is for FixApplicator's own re-join
(\Cref{sec:m-fixapply}). The explanation output is also joined by
\texttt{notebook\_execution\_id}. A coherent run produces at most one
explanation per notebook, and the loader refuses to continue if it finds
two. The fix-application output is joined by \emph{position}: FixApplicator
stamps each result with the line index of the proposal it processed, and
ResultLogger uses that index. Position is needed because
\texttt{notebook\_execution\_id} can legitimately repeat across proposal
lines. The pilot output of \Cref{sec:v-i4} contains a timed-out attempt at
notebook~8 followed by a successful re-run of the same notebook. Joining
by id would not tell the two outcomes apart; joining by position does, and
the two attempts become two rows, as \Cref{subsec:repair-table} requires.

\subsection{Building One Row}
\label{sec:m-resultlogger-row}

\texttt{build\_repair\_attempt\_row()} fills the eighteen columns fixed in
\Cref{subsec:repair-table} from these four inputs. \texttt{failing\_module}
and \texttt{subtype} come from the classification dataset.
\texttt{explanation}, \texttt{llm\_model}, and \texttt{prompt\_strategy}
come from the explanation output; under the orchestrator, a Round-2 row
takes them from the explanation of the newly exposed error
(\Cref{sec:m-orchestrator-record}). The fix and outcome fields come from
the fix-application record when one exists, because that record reflects
what was actually applied and observed. For an abstained proposal, the fix
fields fall back to the proposal's own values, and \texttt{outcome} and
the new-error fields stay \texttt{NULL}. \texttt{rationale} always comes
from the proposal, because FixApplicator makes no repair decision of its
own.

\texttt{run\_id} and \texttt{created\_at} identify the run that produced
the row. When the orchestrator drives every stage in one pass, as in the
final evaluation, it threads one shared \texttt{run\_id} through every
call, so every row of a run carries the same identifier
(\Cref{sec:orchestrator}). Only when a component is run on its own does a
row without an outcome fall back to the repair agent's run identifier.

\subsection{Preserving Structured Evidence}
\label{sec:m-resultlogger-json}

\begin{sloppypar}
Two of the eighteen columns hold a serialized JSON document rather than a
single reduced value. The \texttt{explanation} column holds the complete
explanation object: its \texttt{summary}, \texttt{root\_cause},
\texttt{evidence} list, \texttt{explanation\_confidence}, and
\texttt{limitations}, exactly as validated against
\texttt{schemas/explanation.schema.json}
(\Cref{sec:m-explanation-deviation}). It is not reduced to the
\texttt{summary} field alone. The \texttt{pypi\_evidence} column holds the
complete \texttt{retrieval\_result} object produced by RAGRepairAgent: the
resolved distribution name, every candidate version considered, the
compatibility evidence and its source citation where one was used, and
any warning raised during retrieval. This is more than the three-field
illustration shown in \Cref{subsec:fix-object}. The retrieval object is
serialized whenever the retrieval record exists, which includes a
retrieval-stage abstention: a \texttt{mapping\_unknown} row still carries
a retrieval object whose status and warnings record why package
resolution failed, so every one of the 235 rows of the final evaluation
holds \texttt{pypi\_evidence}. A column is \texttt{NULL} only when the
corresponding upstream field has no value at all, such as the
\texttt{explanation} of a record whose explanation call produced no
response. Neither column ever holds an empty string or an empty JSON
object.
\end{sloppypar}

This preserves detail that a later evaluation might need, without going
back to the original explanation and repair-proposal files, which the
batch that produced them may not have kept. A candidate version that PyPI
offered but the model did not choose, or a compatibility source the
proposal relied on, remains readable from the row that used it, as
evidence and as provenance.

\subsection{From SQLite to CSV}
\label{sec:m-resultlogger-csv}

\begin{sloppypar}
SQLite remains the store of record for evaluation and provenance. The
exporter, \texttt{scripts/\allowbreak export\_repair\_attempts\_csv.py},
adds a CSV export because the existing FAIR Jupyter knowledge-graph
mappings (\texttt{mapping/rml\_mapping/*.rml.ttl}, including
\texttt{executions.rml.ttl}) all read a CSV file through Morph-KGC's
\texttt{ql:CSV} logical source, never a live database connection. The CSV
is a bridge to that existing tooling, not a replacement for the table it
is exported from.
\end{sloppypar}

\begin{sloppypar}
\texttt{repair\_attempts} stores only \texttt{notebook\_execution\_id}
(\Cref{subsec:repair-table}), but the knowledge graph identifies a notebook
by \texttt{notebook\_id}. The exporter resolves \texttt{notebook\_id} for
each row by re-joining \texttt{notebook\_execution\_id} against the same
classification dataset ResultLogger requires, using the loader
FixApplicator already uses for this kind of recovery
(\Cref{sec:m-fixapply}). A \texttt{notebook\_execution\_id} missing from
that dataset raises an error naming the offending
\texttt{repair\_attempts.id}, rather than exporting a row with no notebook
link. By construction this case should never arise, because
\texttt{log\_repair\_attempts()} already required every id it logged to
resolve against the same dataset. The exported columns are \texttt{id} and
\texttt{notebook\_id}, followed by all eighteen \texttt{repair\_attempts}
columns unchanged. \texttt{notebook\_execution\_id} is carried into the CSV
for traceability even though the mapping below does not reference it.
\end{sloppypar}

\subsection{The RML Mapping}
\label{sec:m-resultlogger-rml}

\begin{sloppypar}
\texttt{mapping/rml\_mapping/\allowbreak repair\_attempts.rml.ttl} is modelled
directly on the existing \texttt{executions.rml.ttl}, with the same prefix
block and the same
\texttt{rr:TriplesMap}/\allowbreak\texttt{rr:PredicateObjectMap} structure.
It defines two triples maps over the exported CSV.
\end{sloppypar}

\begin{sloppypar}
The first map gives each repair attempt an IRI,
\url{https://w3id.org/reproduceme/repairattempt_{id}}, typed both
\texttt{repr:RepairAttempt} and \texttt{prov:Activity}. This follows the
vocabulary strategy fixed at design time (\Cref{subsec:kg}): reuse an
existing term where one fits, and add a new \texttt{repr:} term otherwise.
The failing module, the subtype, the explanation, and every field of the
proposed fix and its outcome are attached to this IRI as literal
properties, each with its own new \texttt{repr:} predicate:
\end{sloppypar}

\begin{center}
\small
\begin{tabular}{ll@{\hspace{0.6cm}}ll}
\texttt{failing\_module} & \texttt{repr:failingModule} &
  \texttt{pypi\_evidence} & \texttt{repr:pypiEvidence} \\
\texttt{subtype} & \texttt{repr:subtype} &
  \texttt{outcome} & \texttt{repr:fixOutcome} \\
\texttt{explanation} & \texttt{repr:explanation} &
  \texttt{llm\_model} & \texttt{repr:llmModel} \\
\texttt{action} & \texttt{repr:fixAction} &
  \texttt{prompt\_strategy} & \texttt{repr:promptStrategy} \\
\texttt{install\_name} & \texttt{repr:installName} &
  \texttt{round} & \texttt{repr:round} \\
\texttt{version} & \texttt{repr:fixVersion} &
  \texttt{run\_id} & \texttt{repr:runId} \\
\texttt{command} & \texttt{repr:fixCommand} & & \\
\texttt{rationale} & \texttt{repr:fixRationale} & & \\
\end{tabular}
\end{center}

Two further fields reuse predicates that the existing executions mapping
defines for a cell execution's own exception and message. The new error
type maps to \texttt{repr:exception} and the new error message to
\texttt{repr:msg}. These describe the \emph{post-repair} error, not the
original one. The original failure's exception and message triples already
exist in the graph, written by the executions mapping, so the same kind of
fact is reused for a second observation. Finally, \texttt{created\_at}
maps to \texttt{prov:endedAtTime}, explicitly typed \texttt{xsd:dateTime}.
This is the one place where the mapping departs from the untyped-literal
style of the rest of the graph, because the PROV-O specification gives
\texttt{prov:endedAtTime} an \texttt{xsd:dateTime} range and the
\texttt{created\_at} values are already ISO-8601 strings.

\begin{sloppypar}
A column that is \texttt{NULL} in \texttt{repair\_attempts}, such as an
abstained attempt's \texttt{install\_name} or the \texttt{outcome} of a row
that has none, becomes an empty field in the exported CSV. Morph-KGC then
generates no triple for the corresponding predicate on that row, rather
than an empty-string literal. Missing evidence is therefore absent from
the graph rather than misrepresented as an observed empty value.
\end{sloppypar}

\begin{sloppypar}
The second triples map links the notebook to the repair attempt. Its
subject template, \url{https://w3id.org/reproduceme/notebook_{notebook_id}},
has exactly the same shape as the subject template of
\texttt{notebooks.rml.ttl}, so the two resolve to the same node once
loaded into the same graph. No join and no edit to the existing mapping is
required. The map emits one triple,
\texttt{repr:hadRepairAttempt} \texttt{repr:repairattempt\_\{id\}}, which
attaches the repair history to a notebook resource the graph already
contains.
\end{sloppypar}

This export-and-map chain is the whole of the knowledge-graph side:
ResultLogger writes SQLite, the exporter writes CSV, and Morph-KGC turns
the CSV into RDF through this one mapping file. \Cref{sec:v-i6} reports
the mapping's pilot validation and its subsequent execution, unchanged,
over the complete 235-row table of the final evaluation run.

\subsection{Why No Crosswalk Table Is Needed}
\label{sec:m-resultlogger-identity}

The link above only works if the Docker pipeline's \texttt{notebook\_id}
identifies the same notebook as the \texttt{notebook\_id} of the existing
FAIR Jupyter knowledge graph. That graph was generated from the earlier,
conda-based pipeline (\Cref{sec:pd-why-docker}), not from the Docker
pipeline this thesis builds on. Whether the two pipelines' identifiers
coincide was an open question when the mapping was designed, and it was
settled by checking; \Cref{sec:v-i6} reports the check and its result. The
two pipelines share the same \texttt{notebook}/\texttt{repository} primary
keys for the dataset this thesis operates on. The mapping can therefore
link a repair attempt directly to its existing notebook resource via
\texttt{notebook\_id}, without a separate crosswalk table between the
Docker pipeline and the knowledge graph.

\section{A Worked Example: One Record Through the Pipeline}
\label{sec:m-worked-example}

\begin{sloppypar}
The preceding sections describe each component on its own terms. This
section traces one real record through every stage, to make the data flow
concrete. Every value shown comes from the recorded classification,
explanation, repair-proposal, and fix-application outputs of the pilot runs
reported in \Cref{sec:v-i4,sec:v-i5,sec:v-i7}. The record is
\texttt{notebook\_execution\_id}~174, the \texttt{scipy}/\texttt{cumtrapz}
case already introduced in \Cref{tab:m-compat}.
\end{sloppypar}

\begin{sloppypar}
\paragraph{1. Classification.} The notebook
\texttt{examples/}\allowbreak\texttt{notebooks/}\allowbreak\texttt{Molten\_Salt\_Comparison.ipynb} in
\texttt{zincware/MDSuite} fails with
\texttt{ImportError: cannot import name 'cumtrapz' from}
\texttt{'scipy.integrate'}. Dataset preparation and context extraction
record the failing module as \texttt{scipy}, the refined subtype as
\texttt{wrong\_version}, the scope status as \texttt{usable}, the
root-cause hint as \texttt{version\_or\_api\_}\allowbreak\texttt{incompatibility},
and the classifier's confidence as high.
\end{sloppypar}

\paragraph{2. Explanation.} LLMExplainer, run on this record by the
orchestrator pilot of \Cref{sec:v-i7}, returned a schema-valid explanation
on its first attempt: \emph{summary} "The notebook failed because it tried
to import 'cumtrapz' from 'scipy.integrate', but that name was not
available."; \emph{root cause} "This is likely caused by a package API or
version incompatibility, where the notebook expects a function that is not
present in the installed SciPy package version."; \emph{failing module}
\texttt{scipy}; \emph{confidence} \texttt{medium}, because the installed
SciPy version could not be confirmed from the error alone.
\Cref{sec:m-explanation-example} shows the full six-field record for a
different pilot notebook. Nothing in this explanation is passed on to the
repair step that follows.

\begin{sloppypar}
\paragraph{3. Retrieval.} RAGRepairAgent resolves \texttt{scipy} to itself
(no name mapping needed) and retrieves five Python-3.10-compatible
candidates from the real PyPI Simple API: 1.13.1, 1.13.0, 1.12.0, 1.11.4,
and 1.11.3. Because the subtype is \texttt{wrong\_version}, the
compatibility-evidence registry (\Cref{tab:m-compat}) is also consulted and
resolves \texttt{compatible\_\allowbreak specifier: "<1.14.0"}, citing the SciPy 1.14.0
release notes' "Expired deprecations" entry for \texttt{cumtrapz}.
\end{sloppypar}

\begin{sloppypar}
\paragraph{4. Repair proposal.} Intersecting the candidate list with the
compatibility constraint, the model proposes \texttt{action: "pin\_version"},
\texttt{install\_name: "scipy"}, \texttt{version: "1.13.1"}, with rationale
"cumtrapz was removed in SciPy 1.14.0; 1.13.1 is the newest candidate
version below that boundary, confirmed Python-3.10-compatible." Both schema
and grounding validation pass on the first attempt, and the command is
built deterministically as \texttt{python -m pip install scipy==1.13.1}.
\end{sloppypar}

\paragraph{5. Fix application.} FixApplicator re-joins the record to recover
its repository and commit
(\texttt{eadda45f96874bd6d7eacba89c57daba76db1c7d}), rebuilds the Docker
environment, installs a real \texttt{requirements.txt} and one
\texttt{setup.py} package, then applies the pin. Installation succeeds, and
the notebook is re-executed top to bottom.

\paragraph{6. Observed result.} The original \texttt{cumtrapz} import error
is gone, and execution proceeds further into the notebook. This is direct
evidence that the version pin resolved the targeted incompatibility. The
notebook then fails again on an unrelated
\texttt{ModuleNotFoundError: No module named 'tf\_keras'}. The conservative
\texttt{same\_as\_original\_error} check marks this as \texttt{false}: it
is a different problem, not evidence that the fix did not work. This is
the pattern \Cref{sec:v-key-observation} discusses for the whole pilot:
repairing one dependency can reveal another, and this single attempt has
not made the notebook fully runnable.

In two-round mode, this new error does not simply end the record. The
orchestrator reclassifies \texttt{tf\_keras} as
\texttt{missing\_package}/\texttt{usable}, explains it as a record of its
own (\Cref{sec:m-orchestrator-round2-explain}), and then evaluates its
repair eligibility. The record is eligible, so RAGRepairAgent is called on
it. Retrieval returns \texttt{mapping\_unknown}, because \texttt{tf\_keras}
has no entry in the package-mapping table, and the second round abstains
without a model call (\Cref{sec:v-i7}). The reader still receives an
explanation of the second error; the repair layer, by design, does not
guess at a package it cannot resolve. \Cref{sec:v-case-study} shows an
evaluation-split record where the second round does reach a grounded
repair, with the Round-2 explanation that precedes it.

\paragraph{7. Result logging.} Applying the join logic of
\Cref{sec:m-resultlogger-row} to this record's classification, proposal, and
outcome yields one \texttt{repair\_attempts} row:
\texttt{failing\_module: "scipy"}, \texttt{subtype: "wrong\_version"},
\texttt{action: "pin\_version"}, \texttt{install\_name: "scipy"},
\texttt{version: "1.13.1"}, \texttt{outcome: "still\_failing"},
\texttt{new\_error\_type: "ModuleNotFoundError"},
\texttt{new\_error\_message: "No module named 'tf\_keras'"}, and
\texttt{round: 1}. Because a matching fix-application record exists,
\texttt{run\_id} and \texttt{created\_at} are taken from that FixApplicator
run, as described above. In two-round mode, the Round-2 invocation of
step~6 is logged as well: the \texttt{tf\_keras} abstention becomes a
second row for the same notebook with \texttt{round: 2},
\texttt{action: "none"}, no outcome, and a \texttt{pypi\_evidence} object
whose status is \texttt{mapping\_unknown}
(\Cref{sec:m-orchestrator-record}). This specific Round-1 row was not
produced by a committed ResultLogger run; the mapping's pilot validation
(\Cref{sec:v-i6}) used a different combination of records. The row is
shown here as the mechanical result of the documented join logic applied
to this record's validated outputs, not as a row already in the database.

\begin{sloppypar}
\paragraph{8. Knowledge-graph link.} Exporting the Round-1 row and running
it through \texttt{repair\_attempts.\allowbreak rml.ttl} attaches it to the
existing notebook resource for \texttt{notebook\_id}~428, the notebook's
confirmed identifier (\Cref{sec:m-resultlogger-identity}). The RDF below
shows this Round-1 row only; the Round-2 abstention row would become a
second \texttt{repr:RepairAttempt} resource linked from the same notebook:
\end{sloppypar}

\begin{quote}
\small
\texttt{repr:notebook\_428 repr:hadRepairAttempt repr:repairattempt\_\{id\}} .\\[2pt]
\texttt{repr:repairattempt\_\{id\} a repr:RepairAttempt, prov:Activity} ;\\
\texttt{\ \ repr:failingModule "scipy" ;}\\
\texttt{\ \ repr:fixAction "pin\_version" ;}\\
\texttt{\ \ repr:fixVersion "1.13.1" ;}\\
\texttt{\ \ repr:fixOutcome "still\_failing" ;}\\
\texttt{\ \ repr:exception "ModuleNotFoundError" ;}\\
\texttt{\ \ repr:msg "No module named 'tf\_keras'" .}
\end{quote}

The \texttt{\{id\}} placeholder is left as a template, as in
\Cref{sec:m-resultlogger-rml}, because this row has not been assigned a
primary key by an actual insert. This worked example does not represent a
repair success. It shows what \Cref{sec:v-key-observation} establishes
across the whole pilot: the implemented chain carries a dependency failure
through explanation, grounded repair, re-execution, and structured
logging, and it records that the first repair removed its targeted
dependency error and exposed another dependency error afterwards.

\section{Design Decisions and Trade-offs}
\label{sec:m-decisions}

The sections above explain each decision next to the component it affects.
\Cref{tab:m-tradeoffs} collects the eleven decisions with the largest effect
on how the results in \Cref{chap:validation} should be read. For each
decision it states the alternative that was not taken and the cost of the
choice that was. Several of these costs show up directly in the
evaluation, so the table is worth reading before that chapter. It
introduces no new decision; every row is justified in the section cited.

{\small
\begin{longtable}{|>{\raggedright\arraybackslash}p{2.7cm}|>{\raggedright\arraybackslash}p{2.5cm}|>{\raggedright\arraybackslash}p{3.5cm}|>{\raggedright\arraybackslash}p{3.4cm}|}
  \caption{The design decisions with the largest effect on how the
  evaluation should be interpreted, each stated as decision, alternative,
  reason, and cost.}
  \label{tab:m-tradeoffs}\\
    \hline
    \textbf{Decision} & \textbf{Alternative} & \textbf{Why chosen} & \textbf{Trade-off} \\
    \hline
  \endfirsthead
    \multicolumn{4}{l}{\emph{Table \thetable{}, continued from the previous page.}}\\
    \hline
    \textbf{Decision} & \textbf{Alternative} & \textbf{Why chosen} & \textbf{Trade-off} \\
    \hline
  \endhead
    \hline
    \multicolumn{4}{r}{\emph{continued on the next page}}\\
  \endfoot
    \hline
  \endlastfoot
    Deterministic classification before any LLM call &
      Let a model judge scope and subtype &
      A decision that gates automated repair must be auditable and
      reproducible (\Cref{sec:m-dataset}) &
      Fixed rules miss cases a model might read correctly; 4 of 49 manual
      labels disagree (\Cref{sec:v-final-classifier}) \\
    \hline
    Open, locally deployable model (Gemma-2~9B via Ollama) &
      A proprietary hosted API &
      Open-model requirement NFR1, no hosted-API fee, and the primary
      run's inputs stay on the local machine (\Cref{sec:m-explanation}) &
      CPU-only local serving with high latency, which timed out on 27
      explanation calls in the final Gemma run
      (\Cref{sec:v-final-explanation}) \\
    \hline
    PyPI grounding through an explicit import-to-distribution mapping,
    frozen for the evaluation &
      Free-form repair from the model's memory, or a mapping extended as
      gaps appear &
      Every proposal names a verified package and, for version pins, a
      version from the grounded candidate set; the frozen
      table prevents tuning on the reserved split
      (\Cref{sec:m-repair-retrieval,sec:v-final-manifest}) &
      Abstention whenever the mapping cannot resolve the import:
      162 of 235 invocations (\Cref{sec:v-final-abstention}) \\
    \hline
    Hand-curated compatibility registry for \texttt{wrong\_version} &
      Automated changelog or API-history mining &
      PyPI metadata cannot prove that a symbol exists in a release
      (\Cref{sec:m-repair-compat}) &
      Three entries, scoped to this dataset; unseen incompatibilities
      can only be abstained on \\
    \hline
    Repository dependency files rebuild the baseline environment only;
    they are not repair evidence &
      Hand the declared package list to the repair model &
      Declared files may be incomplete, stale, unpinned, or contain the
      failing version; PyPI evidence is verifiable
      (\Cref{sec:m-repair-nodeps}) &
      Version hints an author did record go unused; a declared pin that
      conflicts with the proposal surfaces only at re-execution
      (\Cref{sec:fw-coverage}) \\
    \hline
    The model never writes the install command &
      The model returns the command string &
      Free-text output must not reach a subprocess; the command is built
      deterministically after validation (\Cref{sec:m-repair}) &
      Only the three structured actions can be expressed \\
    \hline
    Newly exposed Round-2 error explained before its eligibility is checked &
      Explain only errors that will be repaired &
      The reader gets an account of every error encountered; explanation
      and repair stay independent (\Cref{sec:m-orchestrator-round2-explain}) &
      One extra explanation call per new error (50 in the final run) \\
    \hline
    Maximum two repair rounds &
      Open-ended iteration until no error remains &
      A bounded, reproducible evaluation budget
      (\Cref{sec:m-fixapply-round2}) &
      A notebook that exposes another dependency error after Round 2 remains unfixed within the two-round budget.
      (\Cref{sec:v-case-study}) \\
    \hline
    Rebuild the environment on every attempt, with commit pinning where
    metadata is reachable &
      Reuse one long-running container across attempts &
      Clean isolation between attempts, same recipe as the upstream run
      (\Cref{sec:m-fixapply-docker}) &
      Slower, and in the final runs the metadata path was unreachable, so
      default branches were cloned (\Cref{sec:v-final-manifest}) \\
    \hline
    One \texttt{repair\_\allowbreak attempts} row per repair-agent
    invocation, with the full explanation and PyPI evidence kept as JSON &
      Log only applied repairs, or flatten the evidence &
      Every denominator in the evaluation can be reconstructed from the
      table, and abstentions are visible (\Cref{sec:m-resultlogger}) &
      A larger table in which 162 of 235 rows carry no outcome \\
    \hline
    Separate targeted-error-resolution and full-recovery metrics &
      One binary success metric &
      Distinguishes ``the local repair worked'' from ``the whole notebook
      runs'' (\Cref{sec:v-final-targeted}) &
      A more complex evaluation to report and to read \\
\end{longtable}
}
